\documentclass[12pt, a4paper]{article}

% --- PREAMBOLO ---
\usepackage[utf8]{inputenc}
\usepackage[italian]{babel}
\usepackage[T1]{fontenc}
\usepackage{amsmath, amssymb}
\usepackage{geometry} % Gestione margini
\geometry{margin=2.5cm}

% Comando per la linea orizzontale nella copertina
\newcommand{\hr}{\rule{\linewidth}{0.5mm}}

% --- INIZIO DOCUMENTO ---
\begin{document}

% --- COPERTINA ---
\begin{titlepage}
    \centering
    
    {\Large \textbf{Università degli Studi di Padova}} \\
    \vspace{0.3cm}
    {\large Corso di Laurea in Ingegneria Aerospaziale} \\
    \vspace{0.3cm}
    {\normalsize Canale: B} \\
    
    \vspace{2.5cm}
    
    \hr 
    \vspace{0.4cm}
    {\Huge \bfseries Progettino Dinamica del Volo} \\
    \vspace{0.2cm}
    {\large Prima Parte: Descrizione Orbita}
    \vspace{0.4cm}
    \hr
    
    \vspace{3cm}
    
    % Spazio vuoto centrale (al posto dell'immagine)
    \vspace{2cm}
    
    \vfill 
    
    \begin{minipage}{0.45\textwidth}
        \begin{flushleft} \large
            \emph{Autori:}\\
            Alessandro \textbf{Magentini} (Matr. 12345)\\
            Filippo\textbf{Michielan} (Matr. 67890)\\
            Luigi\textbf{Ippolito} (Matr. 12345)\\
            Pietro\textbf{Guglielmini} (Matr. 67890)
        \end{flushleft}
    \end{minipage}
    \begin{minipage}{0.45\textwidth}
        \begin{flushright} \large
            \emph{Docente:} \\
            Prof. Nome \textbf{Cognome}
        \end{flushright}
    \end{minipage}
    
    \vspace{1.5cm}
    {\large \today}
\end{titlepage}

\newpage
\tableofcontents 
\newpage

% --- CONTENUTO ---

\section{Introduzione}
Il primo progettino relativo al corso di Dinamica del Volo, anno accademico 2025/2026, richiede ad un gruppo di massimo quattro studenti di disegnare una missione di intercettazione di un asteroide posizionato su una orbita con inclinazione relativa rispetto al piano equatoriale terrestre.

\section{Analisi del problema}
In questa sezione verranno descritte le equazioni del moto e i parametri orbitali utilizzati per la missione. Un esempio di equazione in blocco separato è il seguente:
\[
\ddot{\vec{r}} + \frac{\mu}{r^3} \vec{r} = 0
\]
\section{Scelta Software specifico}
Per la scelta del programma 
\section{Analisi dei Risultati}
Qui andranno inserite le tabelle dei dati e le descrizioni dei grafici ottenuti.
\section{Descrizione incarichi partecipanti}

\end{document}